\documentclass[a4paper,12pt]{ctexart}
\usepackage{../../teaching-style}

\begin{document}
\pagestyle{empty}

\maintitle{代数}
\begin{center}
  \large 高学段 · 函数、数列、不等式、复数与向量
\end{center}
\vspace{0.3cm}

\chaptertitle{第一章\quad 函数}

\sectiontitle{1.1\quad 函数的概念与三要素}

\knowbox{
函数$y=f(x)$：对于定义域$D$中每一个$x$，有唯一$y$与之对应。\\
三要素：定义域、值域、对应法则。\\
相同函数必须三要素相同。\\
求定义域：分母$\neq0$，偶次根式内$\ge0$，对数真数$>0$等。}

\sectiontitle{1.2\quad 函数的性质}

\knowbox{
\textbf{单调性}：$x_1<x_2\Rightarrow f(x_1)<f(x_2)$为增，反之为减。\\
\textbf{奇偶性}：$f(-x)=f(x)$为偶函数（关于$y$轴对称）；\\
$f(-x)=-f(x)$为奇函数（关于原点对称）。\\
\textbf{周期性}：$f(x+T)=f(x)$，$T$为周期。\\
\textbf{对称性}：若$f(a+x)=f(a-x)$，则关于$x=a$对称。}

\sectiontitle{1.3\quad 基本初等函数}

\knowbox{
\textbf{指数函数}$y=a^x(a>0,a\neq1)$，$x\in\mathbb R$，$y>0$。\\
\textbf{对数函数}$y=\log_a x(a>0,a\neq1)$，$x>0$，$y\in\mathbb R$。\\
\textbf{幂函数}$y=x^\alpha$。\\
\textbf{三角函数}$y=\sin x$，$y=\cos x$，$y=\tan x$（图像与性质）。\\
\textbf{三角恒等变换}：和差角公式、倍角公式、辅助角公式。\\
$\sin(\alpha\pm\beta)=\sin\alpha\cos\beta\pm\cos\alpha\sin\beta$\\
$\cos(\alpha\pm\beta)=\cos\alpha\cos\beta\mp\sin\alpha\sin\beta$\\
$\sin2\alpha=2\sin\alpha\cos\alpha$，$\cos2\alpha=\cos^2\alpha-\sin^2\alpha=2\cos^2\alpha-1=1-2\sin^2\alpha$\\
辅助角：$a\sin x+b\cos x=\sqrt{a^2+b^2}\sin(x+\varphi)$}

\sectiontitle{习题1}

\begin{exercise}
\item 求$f(x)=\dfrac{\sqrt{x+2}}{x-3}$的定义域。
\item 判断奇偶性：$f(x)=x^3$，$f(x)=x^2+1$，$f(x)=x+\dfrac1x$。
\item 已知$f(x+1)=x^2+2x$，求$f(x)$。
\item 化简：$\sin75^\circ$；$\cos105^\circ$；$\sin2\alpha$已知$\sin\alpha=\dfrac35$，$\alpha$为锐角。
\item 求$y=3\sin x+4\cos x$的最大值和最小值。
\end{exercise}

\newpage

\chaptertitle{第二章\quad 数列}

\sectiontitle{2.1\quad 等差数列}

\knowbox{
定义：$a_{n+1}-a_n=d$（常数）。\\
通项：$a_n=a_1+(n-1)d$。\\
前$n$项和：$S_n=\dfrac{n(a_1+a_n)}{2}=na_1+\dfrac{n(n-1)}{2}d$。\\
等差中项：$a,b,c$成等差$\Leftrightarrow2b=a+c$。}

\sectiontitle{2.2\quad 等比数列}

\knowbox{
定义：$\dfrac{a_{n+1}}{a_n}=q$（常数$q\neq0$）。\\
通项：$a_n=a_1q^{n-1}$。\\
前$n$项和：$S_n=\begin{cases}na_1&(q=1)\\\dfrac{a_1(1-q^n)}{1-q}&(q\neq1)\end{cases}$。\\
等比中项：$a,b,c$成等比$\Leftrightarrow b^2=ac$。}

\sectiontitle{2.3\quad 数列求和}

\knowbox{
公式法：等差、等比直接用公式。\\
裂项相消：$\dfrac1{n(n+1)}=\dfrac1n-\dfrac1{n+1}$。\\
错位相减：等差$\times$等比型数列。\\
分组求和：将数列拆成几个可求和的子数列。}

\sectiontitle{2.4\quad 递推数列}

\knowbox{
一阶线性递推：$a_{n+1}=pa_n+q$，构造等比（设$a_{n+1}+\lambda=p(a_n+\lambda)$）。\\
特征根法：二阶线性递推$a_{n+2}=pa_{n+1}+qa_n$，特征方程$x^2-px-q=0$，\\
通项$a_n=A\alpha^n+B\beta^n$（$\alpha,\beta$为特征根）。}

\sectiontitle{习题2}

\begin{exercise}
\item 等差数列$a_n$中，$a_1=3$，$a_5=11$，求$d$、$a_{10}$和$S_{10}$。
\item 等比数列$a_n$中，$a_1=2$，$a_3=18$，求$q$和$S_5$。
\item 求和：$1+\dfrac12+\dfrac14+\cdots+\dfrac1{2^{10}}$。
\item 求和：$\dfrac1{1\times2}+\dfrac1{2\times3}+\cdots+\dfrac1{9\times10}$。
\item 求和：$1\times2+2\times3+3\times4+\cdots+n\times(n+1)$。
\item 递推：$a_1=1$，$a_{n+1}=2a_n+3$，求$a_n$。
\end{exercise}

\newpage

\chaptertitle{第三章\quad 不等式}

\sectiontitle{3.1\quad 不等式的基本性质与证明}

\knowbox{
性质：传递性、可加性（$a>b\Leftrightarrow a+c>b+c$）、可乘性（$c>0$时方向不变，$c<0$时改变）。\\
证明方法：比较法（作差/作商）、综合法、分析法、放缩法。}

\sectiontitle{3.2\quad 均值不等式}

\knowbox{
\textbf{基本不等式}：$\dfrac{a+b}{2}\ge\sqrt{ab}$（$a,b>0$），当$a=b$时取等。\\
推广：$\dfrac{a_1+a_2+\cdots+a_n}{n}\ge\sqrt[n]{a_1a_2\cdots a_n}$。\\
\textbf{柯西不等式}：$(a_1^2+a_2^2)(b_1^2+b_2^2)\ge(a_1b_1+a_2b_2)^2$。\\
一般形式：$(\sum a_i^2)(\sum b_i^2)\ge(\sum a_ib_i)^2$。}

\sectiontitle{3.3\quad 数学归纳法}

\knowbox{
\textbf{第一数学归纳法}：\\
1. 验证$n=n_0$时命题成立。\\
2. 假设$n=k$时成立，证$n=k+1$时也成立。\\
\textbf{第二数学归纳法}：假设$n\le k$时都成立，证$n=k+1$成立。}

\sectiontitle{习题3}

\begin{exercise}
\item 已知$a>0$，$b>0$，求证：$\dfrac{a}{b}+\dfrac{b}{a}\ge2$。
\item $x>0$，求$x+\dfrac4x$的最小值。
\item 已知正数$x,y$满足$x+2y=6$，求$xy$的最大值。
\item 已知$a^2+b^2=1$，$c^2+d^2=1$，求证：$|ac+bd|\le1$。
\item 用数学归纳法证明：$1+3+5+\cdots+(2n-1)=n^2$。
\item 用数学归纳法证明：$2^n>n$（$n\in\mathbb N^*$）。
\end{exercise}

\newpage

\chaptertitle{第四章\quad 复数}

\sectiontitle{4.1\quad 复数的代数形式与运算}

\knowbox{
形如$a+bi$（$a,b\in\mathbb R$，$i^2=-1$）叫复数。\\
$a$为实部，$b$为虚部。\\
加减：$(a+bi)\pm(c+di)=(a\pm c)+(b\pm d)i$\\
乘法：$(a+bi)(c+di)=(ac-bd)+(ad+bc)i$\\
除法：$\dfrac{a+bi}{c+di}=\dfrac{(a+bi)(c-di)}{(c+di)(c-di)}$（分母实数化）。}

\sectiontitle{4.2\quad 复平面与三角形式}

\knowbox{
复平面：横轴为实轴，纵轴为虚轴。\\
模$|z|=\sqrt{a^2+b^2}$，辐角$\theta$满足$\tan\theta=\dfrac ba$。\\
三角形式：$z=r(\cos\theta+i\sin\theta)$。\\
棣莫弗定理：$[r(\cos\theta+i\sin\theta)]^n=r^n(\cos n\theta+i\sin n\theta)$。\\
欧拉公式：$e^{i\theta}=\cos\theta+i\sin\theta$。}

\sectiontitle{习题4}

\begin{exercise}
\item 计算：$(2+3i)+(1-5i)$；$(2+3i)-(1-5i)$。
\item 计算：$(2+3i)(1-5i)$；$\dfrac{2+3i}{1-i}$。
\item 设$z=1+\sqrt3\,i$，求$|z|$、辐角，并化为三角形式。
\item 计算$(1+i)^8$。
\item 解方程：$z^2+2z+5=0$。
\end{exercise}

\newpage

\chaptertitle{第五章\quad 平面向量}

\sectiontitle{5.1\quad 向量的坐标运算}

\knowbox{
坐标表示：$\vec a=(x,y)$。\\
加减：$(x_1,y_1)\pm(x_2,y_2)=(x_1\pm x_2,y_1\pm y_2)$。\\
数乘：$\lambda(x,y)=(\lambda x,\lambda y)$。\\
模：$|\vec a|=\sqrt{x^2+y^2}$。\\
平行：$\vec a\parallel\vec b\Leftrightarrow x_1y_2-y_1x_2=0$。}

\sectiontitle{5.2\quad 数量积}

\knowbox{
$\vec a\cdot\vec b=|\vec a||\vec b|\cos\theta=x_1x_2+y_1y_2$。\\
垂直：$\vec a\perp\vec b\Leftrightarrow\vec a\cdot\vec b=0\Leftrightarrow x_1x_2+y_1y_2=0$。\\
投影：$|\vec a|\cos\theta=\dfrac{\vec a\cdot\vec b}{|\vec b|}$。\\
夹角：$\cos\theta=\dfrac{\vec a\cdot\vec b}{|\vec a||\vec b|}$。}

\sectiontitle{5.3\quad 向量在几何与物理中的应用}

\knowbox{
\textbf{几何}：证平行（$\vec a\parallel\vec b$）、证垂直（$\vec a\cdot\vec b=0$）、\\
求夹角、求距离。\\
\textbf{物理}：力的合成分解、做功$W=\vec F\cdot\vec s$。\\
\textbf{三角形中的向量}：\\
重心：$\overrightarrow{GA}+\overrightarrow{GB}+\overrightarrow{GC}=0$。\\
定比分点：$\overrightarrow{OP}=\dfrac{\overrightarrow{OA}+\lambda\overrightarrow{OB}}{1+\lambda}$。}

\sectiontitle{习题5}

\begin{exercise}
\item $\vec a=(2,-1)$，$\vec b=(1,3)$，求$\vec a+\vec b$、$2\vec a-\vec b$、$\vec a\cdot\vec b$。
\item $\vec a=(1,2)$，$\vec b=(x,4)$，且$\vec a\parallel\vec b$，求$x$。
\item $\vec a=(1,-2)$，$\vec b=(m,4)$，且$\vec a\perp\vec b$，求$m$。
\item $|\vec a|=3$，$|\vec b|=4$，$\vec a$与$\vec b$夹角$60^\circ$，求$\vec a\cdot\vec b$和$|\vec a+\vec b|$。
\item 已知$\triangle ABC$中，$A(1,2)$，$B(3,4)$，$C(5,1)$，求$\overrightarrow{AB}\cdot\overrightarrow{AC}$和$\angle A$。
\end{exercise}

\end{document}
